\documentclass[11pt]{article}
\usepackage[margin=1in]{geometry}
\usepackage[T1]{fontenc}
\usepackage[utf8]{inputenc}
\usepackage{lmodern,amsmath,amssymb,amsthm,mathtools,booktabs,microtype}
\usepackage[colorlinks=true,linkcolor=blue,citecolor=blue,urlcolor=blue]{hyperref}
\hypersetup{pdftitle={Bayesian neural networks compete with small networks in fixed input dimension},pdfauthor={ChatGPT-6 Astra; Claude Opus 5.5; Kaarel Hänni; Dmitry Vaintrob}}
\newtheorem{theorem}{Theorem}[section]
\newtheorem{lemma}[theorem]{Lemma}
\newtheorem{proposition}[theorem]{Proposition}
\newtheorem{corollary}[theorem]{Corollary}
\theoremstyle{definition}\newtheorem{definition}[theorem]{Definition}
\theoremstyle{remark}\newtheorem{remark}[theorem]{Remark}
\newcommand{\E}{\mathbb E}\newcommand{\R}{\mathbb R}\newcommand{\N}{\mathbb N}
\newcommand{\Prob}{\mathbb P}\newcommand{\clip}{\operatorname{clip}}
\newcommand{\KL}{\operatorname{KL}}\newcommand{\Var}{\operatorname{Var}}
\newcommand{\Rad}{\operatorname{Rad}}
\DeclareMathOperator*{\argmin}{arg\,min}
\title{Bayesian neural networks compete with small networks\\in fixed input dimension}
\author{ChatGPT-6 Astra \and Claude Opus 5.5 \and Kaarel H\"anni \and Dmitry Vaintrob\thanks{Authors are listed in alphabetical order.}}
\date{17 November 2026}
\begin{document}
\maketitle
\begin{abstract}
We prove agnostic oracle inequalities for finite Bayesian neural networks
with independent Gaussian weight priors. The data may follow any IID
distribution on a bounded continuous input domain with bounded labels;
no network is assumed to generate the labels. A single probability event
covers every learner width above the comparison network's width and
every Gaussian scaling exponent at or above mean field. The learner
uses its ordinary raw-output posterior mean and a calibrated Gaussian
likelihood. Its risk is bounded by the comparison network's risk with
coefficient one, plus a complexity penalty and a vanishing remainder.
The activation class includes ReLU, GELU, tanh, sigmoid, and softplus,
without requiring $\phi(0)=0$. At fixed input dimension $d$ and depth
$L$, globally Lipschitz activations give fixed-teacher excess
$\widetilde O(n^{-2/[L(d+2)]})$; the sharper bounded-output estimate
for tanh gives
$\widetilde O(n^{-1/\max\{d+2,\,1+dL/2\}})$.
The proof covers functions by their Lipschitz constants, pays Gaussian
norm excursions through a posterior entropy inequality, and controls
raw-output tails. The prior has one fixed dense architecture at each
width, with no mixture over network sizes and no robustness restriction.
The spatial cover incurs a dimensional cost: these bounds complement,
rather than replace, bounds with polynomial dependence on input dimension.
\end{abstract}

\section{Introduction}
A useful approximation guarantee would let a large Bayesian neural
network compete with the best smaller network, even when every small
network has nonzero test error. We establish such a guarantee for dense
Gaussian priors on fully connected networks. The leading coefficient of
the comparison risk is one, so the theorem applies to the population-risk
minimizer of a specified bounded-weight comparison class.

The likelihood must be calibrated. For a learner of hidden width $N$
and prior exponent $\gamma\geq1$, we use Gaussian observation variance
$n\lambda N^{-\gamma}$, where $n$ is the sample count and $\lambda$
is chosen without knowledge of the data distribution. The independent
weight variances are $N^{1-\gamma}/d$ in the first layer,
$N^{-\gamma}$ in subsequent hidden layers, and $N^{-\gamma-1}$ in
the linear readout. These are actual finite-parameter posteriors;
neither width nor prior strength is taken to a limit.

The main argument has three parts. Gaussian matrix tails imply that
large operator norms are exponentially unlikely at the prior's natural
speed. On the complementary event, network outputs have controlled
Lipschitz constants and can be covered by quantizing values on a spatial
grid. An entropy inequality transfers this prior control to the
posterior, while its variational identity compares it with a neighborhood
of any specified teacher. A separate weighted-tail estimate retains the
ordinary, unbounded real output throughout prediction.

Write $H$ for hidden depth and $L=H+1$ for the number of affine maps,
including the readout. At fixed $d$ and $L$, the resulting sample
exponents grow only linearly with depth. This is a statement about the
exponent in a fixed-dimension bound, not a polynomial dependence jointly
on width and depth. The spatial grid introduces an exponential dependence
on input dimension, and constants may also grow rapidly with depth.
No claim is made about efficient posterior computation or gradient descent.

\section{Model and loss}

For input $x\in[-1,1]^d$, put $\bar x=(1,x)$. A width-$N$ network
with $H\geq1$ hidden layers and coordinatewise activation $\phi$ is
\[
 h^1(x)=\phi(W^1\bar x),\qquad
 h^\ell(x)=\phi(W^\ell h^{\ell-1}(x))\ (2\leq\ell\leq H),
 \qquad z_w(x)=v^{\mathsf T}h^H(x).
\]
Here $W^1$ is $N$ by $d+1$, each internal matrix is $N$ by $N$,
and $v\in\mathbb R^N$. Internal layers have no biases. A teacher
uses the same depth and activation, with hidden width $m$ in place
of $N$. Unless stated otherwise its coefficients have absolute value
at most $B\geq1$.

\begin{remark}[Comparison with networks having internal biases]
\label{rem:biased-teacher-compilation}
Suppose $\phi$ is not identically zero. Choose a fixed $a_0$ with
$c=\phi(a_0)\ne0$, and put
\[
 K_\phi=\max\{1,|a_0|,|a_0/c|,1/|c|\}.
\]
Any width-$m$ teacher with a bias at each hidden neuron and at the
readout, with all coefficients and biases bounded by $B\geq1$, can
be represented in the architecture above at width $m+1$, the same
depth, and weight bound $K_\phi B$.

To see this, add one constant neuron at each hidden layer. The first
has preactivation $a_0$ using the constant input; each later one
receives only the previous constant neuron, with weight $a_0/c$.
Every constant neuron therefore outputs $c$. Replace each internal
or readout bias $b$ by an edge of weight $b/c$ from the preceding
constant neuron. All original first-layer biases are already allowed
by the constant input. This preserves every teacher output and gives
the stated coefficient bound. For ReLU take $a_0=c=1$, giving
$K_\phi=1$. For tanh take $a_0=\operatorname{arctanh}(1/2)$ and
$c=1/2$, giving $K_\phi=2$.

This construction concerns only the comparison network. The learner's
Gaussian prior is unchanged, and the resulting bounds apply whenever
$N\geq m+1$. The main theorems themselves do not exclude the zero
activation; only this compilation remark requires a nonzero one.
\end{remark}

The observations $S=((X_i,Y_i))_{i=1}^n$ are IID, with
$X_i\in[-1,1]^d$ and $|Y_i|\leq1$. For any measurable real prediction
function $h$, define
\[
 R(h)=\mathbb E(h(X)-Y)^2,\qquad
 \widehat R_S(h)=\frac1n\sum_{i=1}^n(h(X_i)-Y_i)^2.
\]
The population distribution is otherwise arbitrary. In particular,
no realizability or observation-noise model is imposed on it.

\section{Gaussian norm events and finite function covers}
\label{sec:lipschitz-output-cover}

The first step is a norm event whose probability bound holds at every
learner width. Quantizing a bounded Lipschitz function on a fixed spatial
net then gives a cover independent of that width.

Let $H\geq1$ be the number of hidden layers.  The inputs are
$\bar x=(1,x)$ with $x\in[-1,1]^d$, $d\geq1$, and the readout is
linear.  For $0<\tau\leq1$, the independent centered Gaussian
variances are $\tau^2/d$ in the first layer, $\tau^2/N$ in internal
layers, and $\tau^2/N^2$ in the readout.  Put $\omega=N/\tau^2$.
The canonical exponent $\gamma\geq1$ corresponds to
$\tau^2=N^{1-\gamma}$ and $\omega=N^\gamma$.

\begin{lemma}[An all-width operator-norm event and output cover]
\label{lem:lipschitz-output-cover}
For $K\geq128$, define
\begin{equation}
 G_K=\left\{\|W^1\|_F^2\leq NK,\quad
       \|W^\ell\|_{\rm op}^2\leq K\ (2\leq\ell\leq H),\quad
       N\|v\|_2^2\leq K\right\}.
 \label{eq:lipschitz-norm-event}
\end{equation}
For every $N\geq1$ and every $\tau\in(0,1]$,
\begin{equation}
 \mu_{N,\tau}(G_K^c)\leq(H+1)e^{-\omega K/16}.
 \label{eq:lipschitz-prior-tail}
\end{equation}
If every hidden activation is one-Lipschitz, all outputs on $G_K$
have Lipschitz constant at most
\[
 L_K=K^{(H+1)/2}.
\]
For tanh, the absolute output is at most $A_K=\sqrt K$.
For any one-Lipschitz activation satisfying $\phi(0)=0$, it is at
most $A_K=\sqrt{d+1}\,K^{(H+1)/2}$.

In either case, for every $0<a<1$ there is a deterministic finite
family $\mathcal L_{K,a}$, common to every learner width and every
$\tau\in(0,1]$, with centers bounded by $A_K$, such that all good
outputs are within $a$ in uniform norm of this family and
\begin{align}
 \log|\mathcal L_{K,a}|
 &\leq\left(1+\frac{8\sqrt d\,L_K}{a}\right)^d
                    \log\left(3+\frac{8A_K}{a}\right)\notag\\
 &\leq C_{d,H}K^{d(H+1)/2}a^{-d}
                     \log\frac{C_{d,H}K}{a}.
 \label{eq:lipschitz-output-cover}
\end{align}
The constant $C_{d,H}$ is finite and computable, but is not claimed
to be polynomial in dimension.
\end{lemma}
\begin{proof}
The readout normalization satisfies
$N\|v\|_2^2=(\tau^2/N)\chi_N^2$.
The first-layer normalization satisfies
$\|W^1\|_F^2=(\tau^2/d)\chi_{N(d+1)}^2$.
The Gaussian chi-square exponential moment at parameter $1/4$ gives
\[
 \Pr\{N\|v\|_2^2>K\}
       \leq2^{N/2}e^{-NK/(4\tau^2)}
       \leq e^{-NK/(16\tau^2)},
\]
and
\[
 \Pr\{\|W^1\|_F^2>NK\}
       \leq2^{N(d+1)/2}e^{-dNK/(4\tau^2)}
       \leq e^{-NK/(16\tau^2)}.
\]
The final inequalities use $d\geq1$, $\tau\leq1$, and $K\geq128$.
The standard two-sided sphere-net estimate for each internal matrix is
\[
 \Pr\{\|W^\ell\|_{\rm op}>\sqrt K\}
 \leq2\exp[2N\log9-NK/(8\tau^2)]
 \leq e^{-NK/(16\tau^2)}.
\]
A union bound proves~\eqref{eq:lipschitz-prior-tail}.

Composing Lipschitz maps gives
\[
 \operatorname{Lip}(z_w)
 \leq\|v\|_2\|W^1\|_F
                      \prod_{\ell=2}^H\|W^\ell\|_{\rm op}
 \leq K^{(H+1)/2}.
\]
The last hidden layer has coordinates of magnitude at most one for
tanh, giving $|z_w|\leq\sqrt N\|v\|_2\leq\sqrt K$.
If $\phi(0)=0$, its one-Lipschitz property instead gives
\[
 |z_w(x)|\leq\|v\|_2\|W^1\|_F
           \prod_{\ell=2}^H\|W^\ell\|_{\rm op}\,\|\bar x\|_2
 \leq\sqrt{d+1}\,K^{(H+1)/2}.
\]

Choose a fixed Euclidean net of the cube with radius $a/(4L_K)$,
having at most $(1+8\sqrt d L_K/a)^d$ points.  At each net point
quantize an output value to mesh $a/4$ in $[-A_K,A_K]$, using clipped
endpoints if necessary.  There are at most $3+8A_K/a$ choices at
each point.  Extend every such assignment by the fixed Borel
nearest-net-point map.  The input-net error is at most $a/4$ and the
quantization error at most $a/4$, so the total is strictly less than
$a$.  This proves the first cardinality bound.  Substituting $L_K$
and the relevant $A_K$ gives the second, because $K\geq128$ and
$a<1$.  The construction uses only $d,H,K,a$ and the displayed norm
bounds, so the family is common to all widths and Gaussian scales.
\end{proof}

\begin{lemma}[Uniform prior moments for Lipschitz activations]
\label{lem:lipschitz-prior-moments}
Suppose every hidden activation is one-Lipschitz and vanishes at zero.
Uniformly in $N\geq1$, $0<\tau\leq1$, and $x\in[-1,1]^d$,
\begin{equation}
 \mathbb E z_w(x)=0,\qquad
 \mathbb E z_w(x)^2\leq2\tau^{2(H+1)}/N\leq2,
 \qquad
 \mathbb E z_w(x)^4\leq12\cdot3^H\tau^{4(H+1)}/N^2.
 \label{eq:lipschitz-prior-moments}
\end{equation}
For ReLU, the second-moment bound improves to
$2^{1-H}\tau^{2(H+1)}/N\leq1$.
Consequently, for arbitrary bounded-label IID data and its population
square loss $R$, there is a finite constant $C_H$ such that
\[
 \mathbb E_\mu\widehat R_S(z_w)\leq3,
 \qquad \mathbb E_\mu R(z_w)^2\leq C_H,
\]
uniformly over every sample, every width, and every scale.
The first bound is at most two for ReLU or tanh.
\end{lemma}
\begin{proof}
Put $S_\ell=\|h^\ell(x)\|_2^2/N$ and
$S_0=\|\bar x\|_2^2/d\leq2$.
Conditionally on layer $\ell-1$, each preactivation in layer $\ell$
is an independent centered Gaussian of variance
$\tau^2S_{\ell-1}$.  Since $|\phi(t)|\leq|t|$,
\[
 \mathbb E[S_\ell\mid h^{\ell-1}]\leq\tau^2S_{\ell-1},\qquad
 \mathbb E[S_\ell^2\mid h^{\ell-1}]
       \leq(1+2/N)\tau^4S_{\ell-1}^2
       \leq3\tau^4S_{\ell-1}^2.
\]
The same formula applies at the first layer with the deterministic
quantity $S_0$.  Thus
$\mathbb E S_H\leq2\tau^{2H}$ and
$\mathbb E S_H^2\leq4\cdot3^H\tau^{4H}$.
Conditionally on the final hidden layer, the independent centered
Gaussian readout has variance $\tau^2S_H/N$.  Its exact Gaussian
second and fourth moments prove~\eqref{eq:lipschitz-prior-moments}.
For ReLU, the conditional second-moment recurrence contains the
additional factor $1/2$ at each hidden layer.

The output has prior mean zero, hence its expected empirical square
loss is at most its second moment plus one.  For the second moment
of population risk, Jensen and
$(a-b)^4\leq8(a^4+b^4)$ give the uniform bound
$C_H=8(12\cdot3^H+1)$.  Tanh uses its bounded last-layer activations
to obtain $\mathbb E z_w(x)^2\leq\tau^2/N\leq1$.
\end{proof}

\section{A statistical bound for finite covers}
\label{sec:finite-cover-statistics}

The only uniform statistical estimate needed below is an elementary
finite-class consequence of Bernstein's inequality. The empirical
risk, rather than a worst-case risk bound, appears in its leading term.

\begin{lemma}[Relative deviation for a finite prediction class]
\label{lem:finite-cover-bernstein}
Let $\mathcal F$ be a finite class of measurable real functions, all
bounded in absolute value by $M\geq1$, and let $|Y|\leq1$.
Put $b=(M+1)^2$ and $t=\log(|\mathcal F|/\delta)$.
For $n$ IID observations, with probability at least $1-\delta$,
simultaneously for all $h\in\mathcal F$,
\begin{equation}
 R(h)\leq\widehat R_S(h)
       +3\sqrt{\frac{b\widehat R_S(h)t}{n}}
       +\frac{4bt}{n}.
 \label{eq:finite-cover-bernstein}
\end{equation}
\end{lemma}
\begin{proof}
For one $h$, the squared loss $\ell=(h(X)-Y)^2$ lies in $[0,b]$,
and $\operatorname{Var}(\ell)\leq\mathbb E\ell^2\leq bR(h)$.
The one-sided Bernstein bound and a union bound give
\[
 R(h)-\widehat R_S(h)
 \leq\sqrt{2bR(h)t/n}+2bt/(3n)
\]
simultaneously over the class. Set
$x=(R(h)-\widehat R_S(h))_+$ and $a=2bt/n$.
Then
\[
 x\leq\sqrt{a(\widehat R_S(h)+x)}+a/3
 \leq\sqrt{a\widehat R_S(h)}+x/2+a/2+a/3.
\]
Rearranging gives
$x\leq2\sqrt{2b\widehat R_S(h)t/n}+10bt/(3n)$,
which implies the displayed weaker constants.
\end{proof}

If a network output is uniformly within $a$ of a cover center, then
both population and empirical square-root risks differ by at most $a$,
by Minkowski's inequality. In the proofs below the empirical risk is
bounded by a constant. The cover estimate first bounds the center's
population risk by a constant, after which this square-root comparison
transfers its deviation to the network with a cost proportional to $a$.

\section{Raw oracle bounds for globally Lipschitz activations}
\label{sec:general-lipschitz-raw}

Fix a scalar activation $\phi$ and a constant $A\geq1$ such that
$\operatorname{Lip}(\phi)\leq A$ and $|\phi(0)|\leq A$.
The activation need not be bounded, homogeneous, odd, or zero at the
origin.  In particular the argument covers tanh, ReLU, GELU, sigmoid,
and softplus, with their appropriate fixed Lipschitz constants.

There are $H\geq1$ hidden layers of width $N$, a scalar linear readout,
and a first-layer constant input.  Internal layers have no biases.
Inputs belong to $[-1,1]^d$, labels to $[-1,1]$, and their joint IID
law is arbitrary.  Write $\tau\in(0,1]$ for the actual Gaussian scale:
first-layer, internal, and readout coordinate variances are
$\tau^2/d$, $\tau^2/N$, and $\tau^2/N^2$.
Put $\omega=N/\tau^2$.  The likelihood variance is $n\lambda/\omega$,
so the actual posterior is
\[
 d\Pi=Z^{-1}\exp[-\omega\widehat R_S(z_w)/(2\lambda)]\,d\mu.
\]
Taking $\tau^2=N^{1-\gamma}$ gives the original prefactor-one priors
at every $\gamma\geq1$.  No scale mixture appears in this learner.

\subsection{A teacher cylinder without bounded activations}

Let $g$ be a width-$m$ teacher, with all coefficients at most $B\geq1$
in absolute value.  Define the weighted Gaussian energy
\[
 Q_N(w)=\frac dN\|W^1\|_F^2+
             \sum_{\ell=2}^H\|W^\ell\|_F^2+N\|v\|_2^2.
\]
The Gaussian density has exponent $-\omega Q_N(w)/2$.
To represent the teacher at any $N\geq m$, partition each hidden
layer into clone groups of sizes
$k_j\in\{\lfloor N/m\rfloor,\lceil N/m\rceil\}$.
Copy each teacher first-layer row to its group. For an internal
teacher edge from neuron $i$ to neuron $j$, assign its weight divided
by $k_i$ to every edge between the corresponding clone groups.
Divide each teacher readout weight by its group size in the same way.
The sums entering every clone equal the teacher sums, so this network
computes exactly $g$, for any activation.

The first-layer contribution is at most $d(d+1)B^2$.
Since $k_j/k_i\leq2$, each internal-layer contribution is at most
$2m^2B^2$. Since $N/k_i\leq2m$, the readout contribution is at most
$2m^2B^2$. Thus the clone $w_*$ satisfies
\[
 Q_N(w_*)\leq E:=d(d+1)B^2+2Hm^2B^2.
\]
Every learner neuron belongs to a nonempty clone group, so this
construction does not require $\phi(0)=0$.
Define
\begin{gather*}
 b_0=2\sqrt{2[2\log9+\log(4(H+1))]},\qquad
 T_g=A[3+2mB+B(d+1)+b_0],\\
 P_g=1+(\sqrt2mB+1)A(d+1+H\sqrt m)T_g^{H-1}+T_g^H,\\
 D_g=d+2+(H-1)m,\qquad
 \eta_\rho=\min\{1/2,\rho/(2P_g)\},\qquad c_0=3/8,\\
 \mathcal A_g(\rho)=E/2+D_g\log(1/(c_0\eta_\rho))+\log2.
\end{gather*}
Let $U\geq\|g\|_\infty$ be any specified output bound;
one may always take $U=mBT_g^H$, and for tanh one may take $U=mB$.

\begin{lemma}[A finite-width teacher neighborhood]
\label{lem:general-activation-cylinder}
For every $N\geq m$, $0<\tau\leq1$, and $0<\rho\leq1$,
\[
 -\log\mu\{\|z_w-g\|_\infty<\rho\}
 \leq\omega E/2+ND_g\log(1/(c_0\eta_\rho))+\log2
 \leq\omega\mathcal A_g(\rho).
\]
\end{lemma}
\begin{proof}
Let $w_*$ be the balanced clone.  At each hidden layer let $P$ project
onto the $m$-dimensional subspace of vectors constant on clone groups.
In standardized independent Gaussian coordinates, require absolute
value at most $\eta_\rho$ for all first-layer and readout coordinates,
and for every internal row's projections onto an orthonormal basis of
this group subspace.  There are $ND_g$ independent scalar constraints,
whose joint probability is at least $(c_0\eta_\rho)^{ND_g}$.

Conditioning on these intervals preserves independent centered
unit-subgaussian coordinates in the corresponding orthonormal bases.
Two sphere nets, each of size at most $9^N$, give for an internal
standardized matrix $Z$
\[
 \Pr\{\|Z\|_{\rm op}>b_0\sqrt N\mid\text{cylinder}\}
 \leq2\exp[(2\log9-b_0^2/8)N].
\]
A union bound over internal layers leaves a centrally symmetric event
of probability at least one half the cylinder probability.

Translate this event by $w_*$.  The internal perturbation has operator
norm at most $b_0$, while its restriction to the teacher's group
subspace has operator norm at most $\eta_\rho\sqrt m$.
The clone internal operator norm is at most $\sqrt2mB$.
Teacher hidden coordinates are at most $T_g^\ell$ at layer $\ell$,
because $|\phi(s)|\leq A(1+|s|)$.
For normalized hidden errors $e_\ell=\|h^\ell-h_*^\ell\|_2/\sqrt N$,
Lipschitz propagation gives
\[
 e_1\leq A\eta_\rho(d+1),\qquad
 e_\ell\leq A(\sqrt2mB+b_0)e_{\ell-1}
                +A\eta_\rho\sqrt m\,T_g^{\ell-1}.
\]
Consequently
$e_H\leq\eta_\rho A(d+1+H\sqrt m)T_g^{H-1}$.
The clone readout has $\sqrt N\|v_*\|_2\leq\sqrt2mB$, and its
perturbation has $\sqrt N\|v-v_*\|_2\leq\eta_\rho$.
Thus the output error is at most
$(\sqrt2mB+1)e_H+\eta_\rho T_g^H<\rho$.

For a centered Gaussian measure and a centrally symmetric event,
translation by a vector of squared standardized length $r$ decreases
its probability by at most the factor $e^{-r/2}$; this follows by
pairing opposite points and applying Jensen to the Gaussian density
ratio.  Here $r\leq\omega E$.  This proves the first bound; $N\leq\omega$
and $1\leq\omega$ prove the second.
\end{proof}

\subsection{Width-uniform prior moments and norm tails}

Define
\[
 K(w)=\max\{128,\|W^1\|_F^2/N,
          \|W^2\|_{\rm op}^2,\ldots,\|W^H\|_{\rm op}^2,
          N\|v\|_2^2\}.
\]
Gaussian chi-square tails for the first layer and readout, and two
sphere nets for internal layers, imply for $u\geq128$
\begin{equation}
 \mu\{K>u\}\leq(H+1)e^{-\omega u/16}.
 \label{eq:general-activation-norm-tail}
\end{equation}
For example the internal estimate is
$2\exp[2N\log9-\omega u/8]\leq e^{-\omega u/16}$.
The first-layer term in $K$ is $(\tau^2/(dN))\chi^2_{N(d+1)}$;
the readout term is $(\tau^2/N)\chi^2_N$, so the same weaker bound
applies to them.
Tail integration yields
\begin{equation}
 \log\int e^{\omega K/32}d\mu\leq\omega B_H,
 \qquad B_H=4+\log(H+2).
 \label{eq:general-activation-norm-mgf}
\end{equation}

There are also width-uniform raw-output moments.  Let
$S_\ell=\|h^\ell(x)\|_2^2/N$, and use $S_0=\|\bar x\|_2^2/d\leq2$
for the first conditional Gaussian field.  Conditional on previous
layers, each new field has variance $\tau^2S_{\ell-1}$.
The bounds $|\phi(s)|\leq A(1+|s|)$ and Gaussian second and fourth
moments give
\[
 \mathbb E S_\ell\leq2(4A^2)^\ell,
 \qquad \mathbb E S_\ell^2\leq4(32A^4)^\ell.
\]
For the second assertion, Jensen applied to the average of squared
coordinates and $|\phi(s)|^4\leq8A^4(1+|s|^4)$ suffice; these bounds
also apply to the first layer with the stated $S_0$.
The centered Gaussian readout therefore implies, uniformly in $x$,
\[
 \mathbb E_\mu z_w(x)^2\leq 2(4A^2)^H/N,
 \qquad \mathbb E_\mu z_w(x)^4\leq12(32A^4)^H/N^2.
\]
Put
\[
 R_0=1+2(4A^2)^H,\qquad V_0=8+96(32A^4)^H.
\]
Tonelli and Jensen give, for every sample and every data law under
consideration,
\begin{equation}
 \int\widehat R_S(z_w)d\mu\leq R_0,\qquad
 \int R(z_w)d\mu\leq R_0,\qquad
 \int R(z_w)^2d\mu\leq V_0,
 \qquad Z\geq e^{-\omega R_0/(2\lambda)}.
 \label{eq:general-activation-prior-moments}
\end{equation}

\subsection{A common function cover and raw statistical transfer}

Write $L=H+1$.  On $K(w)\leq k$, Lipschitz propagation gives
\[
 \operatorname{Lip}(z_w)\leq A^H k^{L/2},\qquad
 \|z_w\|_\infty\leq C_{d,H,A}k^{L/2}.
\]
For the amplitude bound, normalized hidden norms satisfy
$a_\ell\leq A(1+\sqrt k\,a_{\ell-1})$ starting from
$a_0\leq\sqrt{d+1}$, and the readout multiplies by at most $\sqrt k$.
Thus this bound includes nonzero $\phi(0)$.
An input grid of mesh proportional to $a/k^{L/2}$, followed by rounding
values to multiples of $a$, provides a deterministic finite cover
$\mathcal C_{k,a}$ in uniform norm, common to all widths and scales,
with centers bounded by $C_{d,H,A}k^{L/2}+1$ and
\begin{equation}
 \log|\mathcal C_{k,a}|
 \leq C_{d,H,A} k^{dL/2}a^{-d}
                  \log\frac{C_{d,H,A}k}{a}
 \qquad(k\geq128,\ 0<a<1).
 \label{eq:general-activation-cover}
\end{equation}
Nearest-grid-point functions are sufficient; the cover centers need
not themselves be networks.

\begin{theorem}[Fixed-dimensional raw oracle inequality]
\label{thm:general-activation-fixed-d}
Fix $d,H,A$ as above and set
\[
 p_{d,H}=\frac{(H+1)(d+2)}2.
\]
There is a finite computable constant $C_{d,H,A}$ such that, whenever
$0<\lambda\leq1$, $0<\delta<1$, $n\geq2$, and
\begin{equation}
 n\geq C_{d,H,A}\lambda^{-p_{d,H}}
             \log^6\!\frac{2n}{\lambda\delta},
 \label{eq:general-activation-budget}
\end{equation}
for every fixed width-$m$, weight-$B$ teacher $g$, with probability
at least $1-\delta$ simultaneously at every $N\geq m$ and every
$0<\tau\leq1$,
\begin{align}
 \int R(z_w)d\Pi
 \leq\;&R(g)+2\lambda\mathcal A_g(\rho)+C_{d,H,A}\lambda\notag\\
 &+(U+1)\sqrt{\frac{2R(g)\log(2/\delta)}n}
       +\frac{2(U+1)^2\log(2/\delta)}{3n},
 \qquad \rho=\frac\lambda{4(U+1)}.
 \label{eq:general-activation-oracle}
\end{align}
The same bound holds for the ordinary raw posterior mean.
In particular it holds on one event for every original exponent
$\gamma\geq1$ and every eligible learner width, without changing the
prior prefactor.  At fixed teacher resources the rate is
$\widetilde O(n^{-2/((H+1)(d+2))})$.
\end{theorem}

\begin{proof}
Take $k_j=128\cdot2^j$ up to the least $J$ with
$k_J\geq(16R_0+128)/\lambda$, and use cover tolerances
$a_j=c_{H,A}\lambda k_j$ with a sufficiently small fixed positive
$c_{H,A}$.  In particular all tolerances are below one.
Apply Lemma~\ref{lem:finite-cover-bernstein} to these finite classes,
dividing failure probability $\delta/2$ among them. Their loss ranges
are at most $C_{d,H,A}k_j^L$.
When $\widehat R_S(z_w)\leq R_0+8$, an approximating center has
empirical risk bounded by a constant depending only on $H,A$.
The finite-cover Bernstein bound therefore makes its population-minus-empirical
risk a sufficiently small multiple of $\lambda k_j$, provided
\[
 n\geq C_{d,H,A}\log^6\!\frac{2n}{\lambda\delta}
    \max_{j\leq J}
       \frac{k_j^{L+dL/2}}{(\lambda k_j)^{d+2}}.
\]
The displayed maximum is at most $C_{H,A}\lambda^{-p_{d,H}}$:
$p_{d,H}\geq d+2$ and $k_j\leq C_{H,A}/\lambda$.
This verifies the condition from~\eqref{eq:general-activation-budget}.
Minkowski transfers population and empirical square-root risks between
a network and its center.  Their risks remain bounded by constants on
this region, so taking $c_{H,A}$ small enough gives
\[
 R(z_w)\leq\widehat R_S(z_w)+\lambda K(w)/16
 \quad\text{when }K(w)\leq k_J,\qquad
                       \widehat R_S(z_w)\leq R_0+8.
\]
Here select a dyadic shell $K(w)\leq k_j\leq2K(w)$ and absorb that
factor into the constants.  This one sample event covers all widths
and all scales, since the finite classes did.

The two omitted regions have small weighted posterior risk.
Using~\eqref{eq:general-activation-prior-moments},
\[
 \int R(z_w)\mathbf1_{\{\widehat R_S>R_0+8\}}d\Pi
 \leq R_0e^{-4\omega/\lambda}.
\]
Cauchy--Schwarz,~\eqref{eq:general-activation-norm-tail}, and the
normalizer bound give
\[
 \int R(z_w)\mathbf1_{\{K>k_J\}}d\Pi
 \leq\sqrt{V_0(H+1)}
     e^{\omega R_0/(2\lambda)-\omega k_J/32}
 \leq\sqrt{V_0(H+1)}e^{-4\omega/\lambda}.
\]
These deterministic estimates are simultaneous in all scales and
widths.  From~\eqref{eq:general-activation-norm-mgf} and posterior
variational optimality, for every comparison law $Q$,
\[
 \int[\widehat R_S+\lambda K/16]d\Pi
 \leq\int\widehat R_SdQ+
          \frac{2\lambda}\omega\operatorname{KL}(Q\Vert\mu)
          +2\lambda B_H.
\]
Choose $Q$ to be the prior conditioned on the teacher cylinder of
radius $\rho$.  Its empirical risk is at most
$\widehat R_S(g)+2\rho(U+1)+\rho^2\leq\widehat R_S(g)+\lambda$,
and Lemma~\ref{lem:general-activation-cylinder} bounds its entropy cost
by $2\lambda\mathcal A_g(\rho)$.
Integrate the deviation bound, add the two tails, and use
$e^{-4/\lambda}\leq\lambda$ to obtain the claimed bound with
$\widehat R_S(g)$ in place of the final teacher-risk terms.
Bernstein's inequality for this fixed teacher supplies those terms
with remaining failure probability $\delta/2$.

Finally the prior fourth moment is finite and the posterior density
relative to the prior is bounded by $1/Z$.  The raw posterior mean
therefore exists, and Jensen proves its risk bound.
\end{proof}

\subsection{Competing with a growing class}

Fix $d,H,A,\delta$, and let $p=p_{d,H}$. Choose the likelihood parameter
of order $\widetilde O(n^{-1/p})$, with a sufficiently large logarithmic
factor to meet \eqref{eq:general-activation-budget}. The same learner
then competes with the population-risk minimizer among teachers whose
width and weight caps are at most $n^\alpha$, whenever
\[
 \alpha<\frac1{4p}=\frac1{2(H+1)(d+2)}.
\]
Its oracle excess is $\widetilde O(n^{-1/p+4\alpha})$.
Indeed the teacher energy is $O(n^{4\alpha})$, the remaining cylinder
cost is $O(n^\alpha\log n)$, and one may take
$U=O(n^{2\alpha(H+1)})$. The minimizer has risk at most one because
zero is an available prediction. Its Bernstein terms therefore vanish
faster than the displayed excess, using $p\geq H+1$ and the stated
range of $\alpha$.

With a fixed weight cap and width cap $n^\alpha$, the sufficient range
and excess are instead
\[
 \alpha<\frac1{2p},\qquad
 \widetilde O(n^{-1/p+2\alpha}).
\]
Compactness of each capped parameter class and continuity of the risk
give a minimizer chosen before sampling. The Gaussian likelihood
calibration depends on the declared resources and confidence level,
not on that minimizer or on the unknown data distribution.

\subsection{A faster bound when averaging clipped outputs}

The likelihood can remain exactly the same while the prediction
convention clips individual network outputs before averaging.  Write
$c(w,x)=\operatorname{clip}_{[-1,1]}z_w(x)$ and
$\widehat f_S^{\rm clip}(x)=\int c(w,x)d\Pi(w)$.
This is generally different from clipping the ordinary raw posterior
mean, so the following refinement specifically concerns the displayed
predictor and individual clipped posterior draws.

\begin{corollary}[Clipped-output oracle bound]
\label{cor:general-activation-clipped}
Theorem~\ref{thm:general-activation-fixed-d} also bounds both
\[
 \int R(c(w,\cdot))d\Pi
 \quad\text{and}\quad R(\widehat f_S^{\rm clip})
\]
under the weaker sample condition
\[
 n\geq C_{d,H,A}\lambda^{-p_{d,H}^{\rm clip}}
          \log^6\!\frac{2n}{\lambda\delta},\qquad
 p_{d,H}^{\rm clip}=\max\{d+2,d(H+1)/2\}.
\]
The actual posterior still uses the raw squared-loss likelihood and
the same single Gaussian prior.
\end{corollary}
\begin{proof}
Clipping preserves the Lipschitz constant and gives output amplitude
one.  The uniform cover construction therefore has the same power
$k^{d(H+1)/2}a^{-d}$ in its logarithmic cardinality, but now a constant
loss range.  Lemma~\ref{lem:finite-cover-bernstein} requires only
\[
 n\geq C_{d,H,A}\log^6\!\frac{2n}{\lambda\delta}
     \max_j\frac{k_j^{d(H+1)/2}}{(\lambda k_j)^{d+2}}.
\]
Since $128\leq k_j\leq C_{H,A}/\lambda$, this maximum is bounded
by a constant times $\lambda^{-\max\{d+2,d(H+1)/2\}}$.
The empirical clipped risk is at most four, so no empirical-risk
truncation is needed.  The resulting common sample event gives
\[
 R(c(w,\cdot))\leq\widehat R_S(c(w,\cdot))+\lambda K(w)/16
             \leq\widehat R_S(z_w)+\lambda K(w)/16
 \qquad(K(w)\leq k_J).
\]
The omitted norm tail is controlled by the same weighted raw-risk
estimate, since clipping cannot increase loss against $|Y|\leq1$.
The entropy and teacher comparisons are unchanged.  Integrating and
then applying Jensen proves the two stated conclusions.
\end{proof}

\section{A sharper bound for tanh}
\label{sec:lipschitz-tanh-raw}

Tanh has an additional useful property: its last hidden layer is bounded
coordinatewise by one. On the same Gaussian norm event, the output has
amplitude $\sqrt K$, even though its Lipschitz constant can be as large
as $K^{(H+1)/2}$. The smaller loss range improves the statistical bound.

Use the preceding architecture with $\phi=\tanh$, and put
\[
 a_{d,H}=1+\frac{d(H+1)}2,\qquad b_d=d+2,
 \qquad s_{d,H}=\max\{a_{d,H},b_d\}.
\]
In the teacher-cylinder definitions take $A=1$ and $U=mB$.
In particular $\rho=\lambda/[4(U+1)]$ and
$\mathcal A_g(\rho)$ is the cost in
Lemma~\ref{lem:general-activation-cylinder}.

\begin{theorem}[A sharper raw oracle bound for tanh]
\label{thm:lipschitz-tanh-raw}
There is a finite computable constant $C_{d,H}$ such that the following
holds. Suppose $n\geq2$, $0<\lambda\leq\epsilon\leq1$,
$0<\delta<1$, and
\begin{equation}
 n\geq C_{d,H}
   \log^6\!\frac{2n}{\lambda\epsilon\delta}\,
   \lambda^{-a_{d,H}}\epsilon^{-(b_d-a_{d,H})_+}.
 \label{eq:lipschitz-tanh-budget}
\end{equation}
For every fixed width-$m$, weight-$B$ teacher $g$ with $B\geq1$,
with probability at least $1-\delta$, simultaneously for all learner
widths $N\geq m$ and all Gaussian scales $0<\tau\leq1$,
\begin{align}
 \int R(z_w)\,\Pi_{N,\tau,S}(dw)
 \leq\;&R(g)+2\lambda\mathcal A_g(\rho)+C_{d,H}\epsilon
 \notag\\
 &+(U+1)\sqrt{\frac{2R(g)\log(2/\delta)}n}
       +\frac{2(U+1)^2\log(2/\delta)}{3n}.
 \label{eq:lipschitz-tanh-oracle}
\end{align}
The ordinary raw-output posterior mean and the clipped prediction
conventions obey the same bound. In particular, one event covers every
original prior exponent $\gamma\geq1$ and every eligible width.

For fixed teacher resources, input dimension, depth, and confidence,
setting $\epsilon=\lambda$ gives excess
\begin{equation}
 \widetilde O\!\left(n^{-1/s_{d,H}}\right).
 \label{eq:lipschitz-tanh-rate}
\end{equation}
At $d=1$, the exponents for two, three, and four hidden layers are
$1/3$, $1/3$, and $2/7$. Joint width and weight caps
$m,B\leq n^\alpha$ are allowed with vanishing oracle excess whenever
\begin{equation}
 \alpha<\frac1{4a_{d,H}}=\frac1{4+2d(H+1)}.
 \label{eq:lipschitz-tanh-growing}
\end{equation}
At a fixed weight cap, the sufficient width exponent is twice as large.
\end{theorem}

\begin{proof}
Use the same continuous norm quantity $K(w)$ and its exponential moment
\eqref{eq:general-activation-norm-mgf}. For tanh, the prior empirical
risk has expectation at most two, and
$\int R(z_w)^2d\mu\leq24$: the last-layer bound gives
$|z_w|^2\leq N\|v\|^2$, whose first and second prior moments are at
most one and three. Thus $Z\geq e^{-\omega/\lambda}$.

Take $k_j=128\cdot2^j$ up to the least $J$ with $k_J\geq160/\lambda$,
and cover tolerances $a_j=c_*(\epsilon+\lambda k_j)$ for a sufficiently
small fixed $c_*>0$. Lemma~\ref{lem:lipschitz-output-cover} supplies
covers of amplitude at most $\sqrt{k_j}$ and logarithmic cardinality
\[
 C_{d,H}k_j^{d(H+1)/2}a_j^{-d}
                 \log(C_{d,H}k_j/a_j).
\]
Apply Lemma~\ref{lem:finite-cover-bernstein} with total failure
probability $\delta/2$. On empirical risk at most ten, the same
square-root-risk transfer as in the preceding proof gives
\[
 R(z_w)\leq\widehat R_S(z_w)+\lambda K(w)/16+C_{d,H}\epsilon
 \quad\text{if } K(w)\leq k_J,
\]
provided the sample count dominates, up to the displayed logarithms,
\begin{equation}
 \max_{128\leq k\leq320/\lambda}
       \frac{k^{a_{d,H}}}{(\epsilon+\lambda k)^{b_d}}.
 \label{eq:lipschitz-tanh-maximization}
\end{equation}
The confidence contribution alone is at most a constant times
$1/(\lambda\epsilon)$ and is also absorbed by
\eqref{eq:lipschitz-tanh-budget}. The proof uses finite covers common
to all widths and scales.

If $a_{d,H}<b_d$, substitute $u=\lambda k/\epsilon$ to bound
\eqref{eq:lipschitz-tanh-maximization} by
\[
 \lambda^{-a_{d,H}}\epsilon^{a_{d,H}-b_d}
      \sup_{u>0}\frac{u^{a_{d,H}}}{(1+u)^{b_d}}
 \leq C_{d,H}\lambda^{-a_{d,H}}
                         \epsilon^{-(b_d-a_{d,H})}.
\]
If $a_{d,H}\geq b_d$, use $\epsilon+\lambda k\geq\lambda k$ and
$k\leq320/\lambda$ to obtain the bound
$C_{d,H}\lambda^{-a_{d,H}}$.
These are the powers in \eqref{eq:lipschitz-tanh-budget}.
There is no minimum-width threshold.

The omitted empirical-risk and norm regions have weighted posterior
risks at most $2e^{-4\omega/\lambda}$ and
$\sqrt{24(H+1)}e^{-4\omega/\lambda}$, respectively, by the same
calculation as in the preceding proof. Integrate the raw deviation,
apply the posterior entropy comparison, and condition the prior on
the teacher cylinder. The tails and the resulting $O_H(\lambda)$
terms are absorbed by $C_{d,H}\epsilon$ because $\lambda\leq\epsilon$.
One fixed-teacher Bernstein event, with failure at most $\delta/2$,
gives the displayed population-risk terms. Jensen proves the statement
for the ordinary posterior mean.

Setting $\epsilon=\lambda$ gives sample power
$a_{d,H}+(b_d-a_{d,H})_+=s_{d,H}$.
For growing teacher caps, choose
$\lambda=n^{-\theta}$ and $\epsilon=n^{-\eta}$ with
\[
 4\alpha<\theta<1/a_{d,H},\qquad
 0<\eta\leq\theta,\qquad
 a_{d,H}\theta+(b_d-a_{d,H})_+\eta<1.
\]
Such choices exist under \eqref{eq:lipschitz-tanh-growing}.
The compact capped class has a population-risk minimizer of risk at
most one, since zero is available. Its energy penalty, logarithmic
cylinder cost, and Bernstein terms all vanish under this calibration.
At a fixed weight cap replace $4\alpha$ by $2\alpha$.
\end{proof}


\section{Scope of the approximation guarantee}
The risk comparison is agnostic: labels need not equal a teacher output,
and their conditional distribution need not match the Gaussian likelihood.
The probability statements require the same distribution for training and
testing. A bounded-weight population-risk minimizer is chosen before the
sample; compactness supplies such a minimizer, and the zero-readout network
shows that its risk is at most one.

The actual prior remains independent Gaussian over all weights of one
dense width-$N$ architecture. Norm events and finite function covers are
proof devices. The posterior is not restricted to those events, and no
prior over architectures is introduced. The main guarantee uses the ordinary raw posterior mean. A separate
corollary gives a faster bound when clipping each network output before
averaging.

The likelihood calibration and comparison weight cap are substantive
assumptions. The theorem does not give the exact finite-sample inequality
with no additive remainder, an unrestricted comparison over arbitrarily
large weights, or a computational procedure for obtaining the posterior.
It also does not eliminate the dimensional cost of approximating general
functions on a continuous domain.

\end{document}
